\scanpage{45}\mspage{44}
donc il faut voir que l'adhérence de $y$ pour ${\unsure{\mathcal{T}|\tau}}'$, et pour $\mathcal{T}$, est pareille (a priori la première serait plus petite). \deleted{Soit en effet $X$ un fermé de $X$, tel que} \inserted{\unsure{Explicitons} en fait}

\uline{Corollaire} Sous \deleted{\unsure{$X$ est une}} \unsure{les|ces} conditions de a) b) c) (sans se restreindre au cas des ultratopologies) si $\unsure{X|A}$ est une partie de $\unsure{U|u|n}$, son adhérence pour $\mathcal{T}'$, et pour $\mathcal{T}$, est la \unsure{$=$}. \deleted{(Idem bien sûr pour une partie de $T$.)}
% slanted
\margin{Cela donne l'adhérence d'une partie quelconque $\deleted{\unsure{X|A}}$ de $S$, pour $\mathcal{T}'$, en termes de $\mathcal{T}$ et de $\mathcal{T}'_0$, \unsure{comme|car} \deleted{$\unsure{\bar{X}^{(\mathcal{T}')}}$}
\[
\bar{X}^{(\mathcal{T}')}=\left(\overline{X\cap U}^{(\mathcal{T})}\right)\cup\overline{(X\cap T)}^{(\mathcal{T}'_0)}
\]}

Soit en effet $X\supset A$, fermé pour $\unsure{T'|\mathcal{T}'}$. alors par les conditions $1^\circ,2^\circ,3^\circ)$ plus haut, $X\supset\overline{X\cap U}^{(\mathcal{T})}$, donc si $X\supset A$, $X\supset\overline{A}^{(\mathcal{T})}$, \deleted{\unsure{a fortiori}} en particulier $\overline{A}^{(\mathcal{T}')}\supset\overline{A}^{(\mathcal{T})}$, OK.

Prouvons e). La première assertion (bas de la page) est une simple \unsure{fonctorialité|formalité|fontorialité}, qui devrait être évidente. \inserted{Les conditions (i) (ii) (iii) pour $\mathcal{T}'$ sont \unsure{dès lors|alors|déjà} connues.} Montrons que la top. de $I$ est \deleted{\unsure{quotient de celle de} $I'$, \unsure{donc}} \inserted{finale pour $J\unsure{\mathbin{\text{---}}|\longrightarrow}I$, $I'\longrightarrow I$.} Cela signifie que pour une partie de $I'$ qui est \uline{$\widetilde{\mathcal{T}}$-saturée} (i.e. de la forme $\varphi^{-1}(A)$), celle-ci est fermée pour $\widetilde{\mathcal{T}}$ si (et seulem\textsuperscript{t} si) elle l'est pour $\overline{\mathcal{T}'}$, \uuline{et} sa trace sur $J'$ est fermée pour $\overline{\mathcal{T}}$. Cela revient à ceci, \deleted{\unsure{une}}
